% Generated from paper/source. Do not edit by hand.
\section{Experimental design}
\subsection{Evidence stages}
The next-phase program was organized so that later stages could not rescue earlier failures by changing the estimand.

\begin{enumerate}
\item \textbf{Comparator audit:} reconstruct the historical result and determine whether the two-channel control is algebraically equivalent.
\item \textbf{Exact-equivalence and mechanism tests:} map the complex operator into a tied real block and test forward and outcome agreement; apply destructive phase interventions only as mechanism probes.
\item \textbf{Power pilot:} estimate the number of worlds required for a minimum practical effect of 0.02 without making an outcome claim.
\item \textbf{Reduced discovery:} train five models on four independent generator families; compute cellwise best-real effects; select one law from six prespecified families.
\item \textbf{Freeze:} store coefficients, thresholds, hashes, family split, and a prohibition on using the candidate for the grand outcome.
\item \textbf{Held-out reduced confirmation:} evaluate the frozen law and architecture effect on four untouched families with new seeds.
\item \textbf{Grand preflight:} open no sealed outcome unless every readiness gate passes.
\end{enumerate}

\subsection{Sample structure}
Discovery trained 200 model instances: four families x two training sizes x five training seeds x five models. Their evaluation generated 14,400 metric records and 2,880 complex-minus-best-real nested effects across 24 worlds and three severities. Held-out evaluation trained 100 instances: four families x one training size x five seeds x five models. It generated 9,600 records and 1,920 nested effects across 32 worlds and three severities.

The hierarchical summary gives equal weight to generator family and treats world as the top-level resampling unit, with training-seed observations nested within world. The held-out analysis uses 20,000 bootstrap resamples over four families, 128 worlds, and 640 family-world-seed observations. A two-sided world-level sign-flip test is reported as exploratory because the reduced profile is outcome-ineligible.

\subsection{Frozen decision thresholds}
A favorable architecture result required a positive aggregate effect whose multiplicity-adjusted interval exceeded +0.02. Law confirmation required all three of R2 at least 0.50, effect-sign accuracy at least 0.80, and MAE at most 0.015. The grand experiment required all readiness gates, not a majority vote. The benchmark-access code therefore returned before execution when any mandatory gate failed.

\subsection{Outcome eligibility}
The discovery and held-out studies were intentionally useful but insufficient substitutes for the full program. Their protocol deviations were written into each result artifact: omitted training sizes, incomplete real envelope, one-trial search, incomplete ShiftGauntlet, and no authority to support the strongest outcome category or QN-GRAND-001. We report their numerical results because they are direct evidence about the candidate claim; we do not relabel them as the preregistered grand confirmation.

\subsection{Statistical estimands}
The principal descriptive estimand is mean complex top-1 minus cellwise best-real top-1. We also report the median, trimmed mean, worst decile, and probability of superiority. Exact-real checks use maximum absolute metric differences. The law uses 12 family-by-severity aggregate cells in discovery and 12 in held-out evaluation. All result directories, analysis code, claim text, and figure inputs are distributed with the repository.
